\documentclass[12pt,a4paper]{article}

\usepackage[utf8]{inputenc}
\usepackage[T1]{fontenc}
\usepackage[polish]{babel}
\usepackage{geometry}
\usepackage{setspace}
\usepackage{hyperref}
\usepackage{graphicx}
\usepackage{listings}
\usepackage{xcolor}
\usepackage{float}

\graphicspath{{screenshots/}}

\geometry{margin=2.5cm}
\onehalfspacing

\definecolor{codebg}{HTML}{ffffff}
\definecolor{codetext}{HTML}{1f2328}
\definecolor{keyword}{HTML}{6f42c1}
\definecolor{string}{HTML}{0a7d20}
\definecolor{comment}{HTML}{6a737d}
\definecolor{annotation}{HTML}{b08000}
\definecolor{framecolor}{HTML}{d0d7de}

\lstdefinestyle{java}{
  language=Java,
  basicstyle=\ttfamily\footnotesize\color{codetext},
  backgroundcolor=\color{codebg},
  commentstyle=\color{comment}\itshape,
  keywordstyle=\color{keyword}\bfseries,
  stringstyle=\color{string},
  numberstyle=\tiny\color{comment},
  numbers=left,
  numbersep=8pt,
  breaklines=true,
  frame=single,
  rulecolor=\color{framecolor},
  tabsize=4,
  showstringspaces=false,
  morekeywords={var, record, yield},
  literate={ą}{{\k{a}}}1 {ć}{{\'c}}1 {ę}{{\k{e}}}1 {ł}{{\l{}}}1 {ń}{{\'n}}1 {ó}{{\'o}}1 {ś}{{\'s}}1 {ź}{{\'z}}1 {ż}{{\.z}}1
            {Ą}{{\k{A}}}1 {Ć}{{\'C}}1 {Ę}{{\k{E}}}1 {Ł}{{\L{}}}1 {Ń}{{\'N}}1 {Ó}{{\'O}}1 {Ś}{{\'S}}1 {Ź}{{\'Z}}1 {Ż}{{\.Z}}1,
}

\lstdefinestyle{typescript}{
  language=Java,
  basicstyle=\ttfamily\footnotesize\color{codetext},
  backgroundcolor=\color{codebg},
  commentstyle=\color{comment}\itshape,
  keywordstyle=\color{keyword}\bfseries,
  stringstyle=\color{string},
  numberstyle=\tiny\color{comment},
  numbers=left,
  numbersep=8pt,
  breaklines=true,
  frame=single,
  rulecolor=\color{framecolor},
  tabsize=2,
  showstringspaces=false,
  morekeywords={const, let, async, await, import, from, export, default, type, interface, function, return, useState, useEffect, useMemo, Promise},
  morecomment=[l]{//},
  morestring=[b]',
  morestring=[b]`,
}

\lstdefinestyle{yaml}{
  basicstyle=\ttfamily\footnotesize,
  backgroundcolor=\color{codebg},
  commentstyle=\color{comment}\itshape,
  keywordstyle=\color{keyword}\bfseries,
  numberstyle=\tiny\color{comment},
  numbers=left,
  numbersep=8pt,
  breaklines=true,
  frame=single,
  rulecolor=\color{codebg},
  tabsize=2,
  showstringspaces=false,
  morecomment=[l]{\#},
}

\title{\textbf{Sprawozdanie z~Projektu:\\ Trading Journal}}
\author{
  Daniel Salawa, Igor Świerczek \\
  ITE, 1 rok, 2 stopień
}
\date{}

\begin{document}

\maketitle

% ============================================================
\section{Cel projektu}
% ============================================================

Celem projektu było zaprojektowanie i~zaimplementowanie zaawansowanej platformy webowej \textbf{Trading Journal}, wspomagającej proces decyzyjny i~analityczny inwestorów giełdowych. Aplikacja miała zapewnić pełną kontrolę nad historią transakcji, umożliwić analizę psychologiczną tradera, dokumentowanie strategii inwestycyjnych oraz wizualizację wyników handlowych. Całość miała zostać zrealizowana w~architekturze klient-serwer z~wykorzystaniem nowoczesnych technologii webowych oraz zewnętrznych integracji finansowych.

\subsection*{Podział obowiązków}
\begin{itemize}
  \item \textbf{Daniel Salawa}, backend i~infrastruktura: logika biznesowa (Spring Boot), baza danych (PostgreSQL + Flyway), integracje z~API zewnętrznymi, konteneryzacja (Docker), wdrożenie (Railway).
  \item \textbf{Igor Świerczek}, frontend i~bezpieczeństwo: interfejs użytkownika (React + TypeScript), system autentyfikacji JWT, zabezpieczenia (CORS, BCrypt, walidacja dostępu), wizualizacje danych (Recharts).
\end{itemize}

% ============================================================
\section{Użyte technologie}
% ============================================================

\begin{table}[H]
\centering
\begin{tabular}{|l|l|}
\hline
\textbf{Warstwa / Obszar} & \textbf{Technologia} \\ \hline
Backend & Java 17, Spring Boot 3, Spring Security, Lombok \\ \hline
Baza danych & PostgreSQL 16, Flyway (13 migracji) \\ \hline
Frontend & React 18, TypeScript, Vite, Tailwind CSS \\ \hline
Wykresy & Recharts, TradingView Widget \\ \hline
Autentyfikacja & JWT (HMAC), BCrypt \\ \hline
API zewnętrzne & Finnhub (akcje), ExchangeRate-API (forex), ForexFactory (makro) \\ \hline
Konteneryzacja & Docker Compose (PostgreSQL, backend, frontend + Nginx) \\ \hline
Wdrożenie & Railway (backend), Vercel (frontend) \\ \hline
Kontrola wersji & Git, GitHub \\ \hline
\end{tabular}
\caption{Zestawienie użytych technologii}
\end{table}

% ============================================================
\section{Zakładane funkcjonalności}
% ============================================================

\begin{enumerate}
  \item \textbf{Rejestracja i~logowanie}, bezpieczna autentyfikacja oparta na tokenach JWT z~hashowaniem haseł BCrypt oraz resetem hasła.
  \item \textbf{Dashboard}, przegląd konta: saldo, łączny P\&L, win rate, tabela transakcji z~sortowaniem i~filtrami.
  \item \textbf{Zarządzanie transakcjami (Trade)}, tworzenie, edycja i~usuwanie transakcji z~parametrami: cena wejścia/wyjścia, kierunek (LONG/SHORT), stop loss, rozmiar pozycji, emocje, ocena, notatki, tagi, załączniki graficzne.
  \item \textbf{Playbook (Strategie)}, dokumentowanie strategii handlowych w~formacie Markdown z~dynamicznymi checklistami (JSON), automatycznie pojawiającymi się w~formularzu transakcji.
  \item \textbf{Dziennik psychologiczny}, codzienne wpisy z~oceną nastroju (1--5), poziomu energii, notatek i~wniosków.
  \item \textbf{Analytics}, zaawansowane metryki (Win Rate, Profit Factor, Expectancy, Sharpe Ratio, Max Drawdown) i~wizualizacje (krzywa equity, rozkład wg dnia tygodnia, proporcja win/loss, wydajność wg strategii).
  \item \textbf{Kalendarz tradingowy}, widok kalendarza łączący transakcje z~wydarzeniami makroekonomicznymi (ForexFactory).
  \item \textbf{Long Term (Planowanie strategiczne)}, moduł organizujący myślenie w~kategoriach: THEME, WATCHLIST, GOAL, REVIEW.
  \item \textbf{Chart}, integracja z~widgetem TradingView do analizy technicznej.
  \item \textbf{Panel administracyjny}, zarządzanie użytkownikami, statystyki systemowe, konfiguracja klucz-wartość.
  \item \textbf{Konteneryzacja i~wdrożenie}, Docker Compose dla środowiska lokalnego, deployment na Railway/Vercel.
\end{enumerate}

% ============================================================
\section{Fragmenty kodu z~omówieniem}
% ============================================================

% ---- Rysunek 1 ----
\begin{figure}[H]
\begin{lstlisting}[style=java, caption={}, label={}]
public Trade toEntity(User user, TradeRequest request,
                      Playbook playbook) {
    return Trade.builder()
            .user(user)
            .ticker(request.getTicker().toUpperCase())
            .direction(request.getDirection())
            .entryPrice(request.getEntryPrice())
            .exitPrice(request.getExitPrice())
            .positionSize(request.getPositionSize())
            .stopLoss(request.getStopLoss())
            .playbook(playbook)
            .notes(request.getNotes())
            .tags(request.getTags())
            .emotionBefore(request.getEmotionBefore())
            .emotionAfter(request.getEmotionAfter())
            .preTradeChecklist(request.getPreTradeChecklist())
            .postTradeReview(request.getPostTradeReview())
            .build();
}

private Double calculatePnl(Trade trade) {
    if (trade.getExitPrice() == null) return null;
    double multiplier =
        trade.getDirection() == TradeDirection.LONG ? 1 : -1;
    return (trade.getExitPrice() - trade.getEntryPrice())
           * trade.getPositionSize() * multiplier;
}
\end{lstlisting}
\caption{TradeMapper, mapowanie DTO na encję oraz obliczanie P\&L. Metoda \texttt{calculatePnl} uwzględnia kierunek transakcji: dla pozycji SHORT stosuje odwrócony mnożnik, dzięki czemu zysk jest poprawnie obliczany niezależnie od kierunku.}
\label{fig:trade-mapper}
\end{figure}

% ---- Rysunek 2 ----
\begin{figure}[H]
\begin{lstlisting}[style=java, caption={}, label={}]
@Override
protected void doFilterInternal(
        HttpServletRequest request,
        HttpServletResponse response,
        FilterChain filterChain) throws ServletException,
                                       IOException {
    String header = request.getHeader("Authorization");
    if (StringUtils.hasText(header)
            && header.startsWith("Bearer ")) {
        String token = header.substring(7);

        if (jwtUtil.isTokenValid(token)
                && SecurityContextHolder.getContext()
                   .getAuthentication() == null) {
            Claims claims = jwtUtil.parseClaims(token);
            String email = claims.getSubject();
            String roleStr = claims.get("role", String.class);
            Role role = roleStr != null
                ? Role.valueOf(roleStr) : Role.USER;

            UsernamePasswordAuthenticationToken auth =
                new UsernamePasswordAuthenticationToken(
                    email, null,
                    List.of(new SimpleGrantedAuthority(
                        role.authority())));
            auth.setDetails(
                new WebAuthenticationDetailsSource()
                    .buildDetails(request));
            SecurityContextHolder.getContext()
                .setAuthentication(auth);
        }
    }
    filterChain.doFilter(request, response);
}
\end{lstlisting}
\caption{JwtAuthFilter, filtr bezpieczeństwa Spring Security. Przy każdym żądaniu HTTP wyciąga token JWT z~nagłówka \texttt{Authorization}, weryfikuje podpis HMAC, odczytuje email i~rolę użytkownika, a~następnie ustawia kontekst bezpieczeństwa, aby dalsze warstwy aplikacji znały tożsamość użytkownika.}
\label{fig:jwt-filter}
\end{figure}

% ---- Rysunek 3 ----
\begin{figure}[H]
\begin{lstlisting}[style=java, caption={}, label={}]
@Bean
public SecurityFilterChain filterChain(HttpSecurity http,
        JwtUtil jwtUtil) throws Exception {
    http
        .csrf(AbstractHttpConfigurer::disable)
        .cors(Customizer.withDefaults())
        .sessionManagement(s -> s.sessionCreationPolicy(
            SessionCreationPolicy.STATELESS))
        .authorizeHttpRequests(auth -> auth
            .requestMatchers("/api/auth/**").permitAll()
            .requestMatchers("/actuator/health").permitAll()
            .requestMatchers(HttpMethod.OPTIONS, "/**")
                .permitAll()
            .requestMatchers("/api/admin/**").hasRole("ADMIN")
            .anyRequest().authenticated()
        )
        .addFilterBefore(new JwtAuthFilter(jwtUtil),
            UsernamePasswordAuthenticationFilter.class);
    return http.build();
}
\end{lstlisting}
\caption{SecurityConfig, konfiguracja łańcucha filtrów Spring Security. API jest bezstanowe (STATELESS, brak sesji). Endpointy autentyfikacji są publiczne, panel admina wymaga roli ADMIN, a~pozostałe zasoby wymagają uwierzytelnienia tokenem JWT.}
\label{fig:security-config}
\end{figure}

% ---- Rysunek 4 ----
\begin{figure}[H]
\begin{lstlisting}[style=java, caption={}, label={}]
@Transactional
public List<AttachmentDto> uploadAttachments(
        String email, UUID tradeId,
        List<MultipartFile> files) throws IOException {
    ensureOwned(email, tradeId);
    var allowed = Set.of("image/jpeg", "image/png",
                         "image/webp", "image/gif");
    List<TradeAttachment> saved = new ArrayList<>();

    for (MultipartFile f : files) {
        if (f.isEmpty()) continue;
        String url = fileUploadService
            .saveFile(f, allowed);
        var att = TradeAttachment.builder()
                .tradeId(tradeId)
                .filename(f.getOriginalFilename())
                .contentType(f.getContentType())
                .size(f.getSize())
                .url(url)
                .build();
        saved.add(attachmentRepository.save(att));
    }
    return saved.stream()
        .map(this::toAttachmentDto).toList();
}
\end{lstlisting}
\caption{TradeService, przesyłanie załączników graficznych. Metoda weryfikuje przynależność transakcji do użytkownika (\texttt{ensureOwned}), waliduje typ MIME pliku (tylko obrazy), zapisuje plik na dysku z~unikalną nazwą UUID i~tworzy rekord w~bazie danych.}
\label{fig:upload-attachments}
\end{figure}

% ---- Rysunek 5 ----
\begin{figure}[H]
\begin{lstlisting}[style=typescript, caption={}, label={}]
// Interceptor Axios, automatyczne dolaczanie tokenu JWT
api.interceptors.request.use((config) => {
  const token = localStorage.getItem('jwt');
  if (token) {
    config.headers.Authorization = `Bearer ${token}`;
  }
  return config;
});

// Auto-logout przy 401, retry przy 502/503/504
api.interceptors.response.use(
  (res) => res,
  async (err: AxiosError) => {
    if (err?.response?.status === 401) {
      localStorage.removeItem('jwt');
      location.href = '/';
    }
    if (isRetryableStartupError(err)) {
      const config = err.config as RetryableRequestConfig;
      config._retryCount = (config._retryCount ?? 0) + 1;
      if (config._retryCount <= MAX_STARTUP_RETRIES) {
        await wait(config._retryCount * 700);
        return api.request(config);
      }
    }
    return Promise.reject(err);
  },
);
\end{lstlisting}
\caption{Klient API (frontend), interceptory Axios. Pierwszy interceptor automatycznie dołącza token JWT do każdego żądania. Drugi obsługuje błędy: przy statusie 401 wylogowuje użytkownika, a~przy błędach 502/503/504 automatycznie ponawia żądanie (do 3 prób z~rosnącym opóźnieniem).}
\label{fig:axios-interceptors}
\end{figure}

% ---- Rysunek 6 ----
\begin{figure}[H]
\begin{lstlisting}[style=typescript, caption={}, label={}]
const stats = useMemo(() => {
  const initialBalance = me?.initialBalance ?? 10000;
  if (!trades.length)
    return { totalPL: 0, winRate: 0, totalTrades: 0,
             wins: 0, losses: 0,
             currentBalance: initialBalance };

  let totalPL = 0, wins = 0, losses = 0;
  for (const tr of trades) {
    const pl = calculatePL(tr);
    totalPL += pl;
    if (pl > 0) wins++;
    else if (pl < 0) losses++;
  }

  return {
    totalPL,
    winRate: Math.round((wins / trades.length) * 100),
    totalTrades: trades.length, wins, losses,
    currentBalance: initialBalance + totalPL,
  };
}, [trades, me]);
\end{lstlisting}
\caption{Dashboard, obliczanie statystyk w~czasie rzeczywistym przy użyciu hooka \texttt{useMemo}. Statystyki (saldo, P\&L, win rate) są przeliczane tylko wtedy, gdy zmieni się lista transakcji, co minimalizuje zbędne renderowanie komponentów React.}
\label{fig:dashboard-stats}
\end{figure}

% ---- Rysunek 7 ----
\begin{figure}[H]
\begin{lstlisting}[style=yaml, caption={}, label={}]
services:
  db:
    image: postgres:16
    environment:
      POSTGRES_DB: journal_db
      POSTGRES_USER: postgres
      POSTGRES_PASSWORD: postgres
    ports:
      - "5443:5432"
    volumes:
      - pg_data:/var/lib/postgresql/data

  backend:
    build: ./backend
    ports:
      - "8080:8080"
    environment:
      SPRING_PROFILES_ACTIVE: dev
      SPRING_DATASOURCE_URL:
        jdbc:postgresql://db:5432/journal_db
    depends_on:
      - db

  frontend:
    build: ./frontend
    ports:
      - "5173:8080"
    environment:
      BACKEND_URL: http://backend:8080
    depends_on:
      - backend

volumes:
  pg_data:
\end{lstlisting}
\caption{Docker Compose, definicja trzech kontenerów: baza PostgreSQL 16 z~wolumenem persystentnym, backend Spring Boot z~profilem deweloperskim oraz frontend (Nginx) z~reverse proxy do backendu. Zależności (\texttt{depends\_on}) zapewniają poprawną kolejność uruchamiania.}
\label{fig:docker-compose}
\end{figure}

% ============================================================
\section{Zrzuty ekranowe z~działającej aplikacji}
% ============================================================

W~tej sekcji prezentujemy widoki z~działającej aplikacji. Pod każdym zrzutem znajduje się opis tego, co prezentuje dany widok i~jakie elementy interfejsu są na nim widoczne.

% Wstawia zrzut ekranu z~katalogu screenshots/ (format poziomy)
\newcommand{\screenshotPlaceholder}[2]{%
  \begin{figure}[H]
    \centering
    \includegraphics[width=0.95\textwidth]{#1.png}
    \caption{#2}
    \label{fig:#1}
  \end{figure}%
}

% Wstawia zrzut ekranu w~formacie pionowym (modale, waskie widoki)
\newcommand{\screenshotPlaceholderTall}[2]{%
  \begin{figure}[H]
    \centering
    \includegraphics[width=0.70\textwidth]{#1.png}
    \caption{#2}
    \label{fig:#1}
  \end{figure}%
}

% --- 1. Logowanie / rejestracja ---
\subsection*{Logowanie i~rejestracja}
Strona startowa aplikacji prezentuje wycentrowany formularz autentyfikacji z~przełącznikiem między trybem logowania (Sign~In) a~rejestracji (Sign~Up). W~trybie logowania pole hasła wyposażone jest w~przycisk \texttt{Show/Hide} ujawniający wpisany ciąg, a~pod polem hasła znajduje się link \emph{Forgot~password?} prowadzący do~procedury resetu. W~trybie rejestracji formularz rozszerzany jest o~pole \texttt{Username} (autouzupełnianie HTML5: \texttt{username}) i~automatycznie zmienia atrybut \texttt{autocomplete} dla hasła na \texttt{new-password}. Po~poprawnym uwierzytelnieniu backend zwraca token JWT, który zapisywany jest w~\texttt{localStorage} i~następnie automatycznie dołączany przez interceptor Axios do~każdego kolejnego żądania (Rys.~\ref{fig:axios-interceptors}). Komunikaty błędów (nieprawidłowe dane, zablokowane konto) są pokazywane w~osobnym banerze z~animacją wjazdu.
\screenshotPlaceholder{login}{Ekran logowania, formularz uwierzytelnienia opartego na~tokenie JWT}

% --- 2. Dashboard ---
\subsection*{Dashboard, widok główny}
Dashboard jest ekranem startowym po~zalogowaniu. U~góry strony znajduje się powitanie z~nazwą użytkownika oraz siatka kart statystycznych: \textbf{Account~Balance} (saldo początkowe + skumulowany P\&L, z~przyciskiem \emph{Edit} otwierającym modal edycji salda), \textbf{Total~P\&L} (z~ikoną kierunku trendu), \textbf{Win~Rate} (procentowo, kolorowane zielono/czerwono w~zależności od~poziomu) oraz karta-CTA \textbf{Add~New~Trade} otwierająca modal dodawania transakcji. Pod~kartami znajduje się pasek filtrów (wyszukiwanie ticker/notatki, kierunek LONG/SHORT, wynik win/loss/break-even) oraz tabela ostatnich transakcji z~klikalnymi nagłówkami kolumn (sortowanie wzn./mal. po~dacie, tickerze, wielkości pozycji i~P\&L). Każdy wiersz tabeli pokazuje datę, ticker, kolorowany badge LONG/SHORT, ceny wejścia i~wyjścia, P\&L (kolorowane), badge wyniku oraz licznik załączników. Wszystkie statystyki są przeliczane po~stronie klienta przy użyciu hooka \texttt{useMemo} (Rys.~\ref{fig:dashboard-stats}), co~zapewnia natychmiastową reakcję interfejsu na~każdą zmianę listy transakcji.
\screenshotPlaceholder{dashboard}{Dashboard, karty statystyczne i~tabela transakcji z~filtrami}

% --- 3. Trade Modal ---
\subsection*{Modal dodawania i~edycji transakcji}
Modal transakcji jest najbogatszym formularzem w~aplikacji, dostępnym z~przycisku \emph{Add~New~Trade} oraz z~każdego wiersza tabeli (otwiera się w~trybie edycji). Formularz dzieli się na~kilka logicznych sekcji: \textbf{Identyfikacja} (ticker, data otwarcia/zamknięcia), \textbf{Pozycja} (kierunek LONG/SHORT, cena wejścia/wyjścia, stop loss, wielkość pozycji), \textbf{Strategia} (lista rozwijana z~playbookami użytkownika, dynamicznie wstrzykująca checklistę pre-trade), \textbf{Psychologia} (ocena emocji przed i~po~transakcji, ocena 1--10, tagi setupu) oraz \textbf{Załączniki} (drag~\&~drop dla zrzutów wykresów, przechowywanych w~katalogu uploads na~backendzie z~zabezpieczeniem przed path traversal). Modal posiada również automatyczny \emph{position~sizing} obliczający rekomendowaną wielkość pozycji na~podstawie ryzyka procentowego do~saldo konta i~odległości stop lossa, oraz przycisk \emph{Fetch~current~price} pobierający aktualny kurs z~Finnhub/ExchangeRate-API.
\screenshotPlaceholderTall{trade-modal}{Modal dodawania transakcji, pełny formularz z~sekcjami identyfikacji, pozycji i~psychologii}

% --- 4. Analytics ---
\subsection*{Moduł Analytics, zaawansowane metryki tradingowe}
Strona Analytics jest centralnym punktem oceny skuteczności tradera. Na~górze wyświetlana jest siatka dziewięciu kart z~kluczowymi metrykami obliczanymi w~module \texttt{analyticsCalculations.ts}: \textbf{Total~Trades}, \textbf{Win~Rate}, \textbf{Total~P/L}, \textbf{Profit~Factor} (suma wygranych podzielona przez sumę przegranych), \textbf{Avg~Win}, \textbf{Avg~Loss}, \textbf{Expectancy} (oczekiwana wartość pojedynczej transakcji), \textbf{Max~Drawdown} (maksymalny spadek od~szczytu krzywej kapitału, w~jednostkach i~procentowo) oraz \textbf{Sharpe~Ratio} (stosunek średniego zwrotu do~odchylenia standardowego). Pod~metrykami znajduje się szeroki wykres \textbf{Equity~Curve} (Recharts \texttt{LineChart}) pokazujący ewolucję salda po~każdej transakcji. Niżej dwa wykresy kolumnowe/kołowe: \textbf{Trades~by~Day~of~Week} z~rozbiciem na~wygrane i~przegrane oraz \textbf{Win/Loss~Distribution} jako wykres kołowy. Sekcja \textbf{Performance~by~Strategy} prezentuje karty per playbook z~win rate, ilością transakcji i~paskiem postępu. Na~końcu tabela \textbf{Performance~by~Tag/Setup} agreguje wyniki według tagów setupu.
\screenshotPlaceholder{analytics-metrics}{Analytics, siatka kart z~metrykami (Win Rate, Profit Factor, Sharpe, Max DD)}
\screenshotPlaceholder{analytics-equity}{Analytics, krzywa kapitału (Equity Curve) i~wykresy rozkładu}
\screenshotPlaceholder{analytics-strategy}{Analytics, sekcja \emph{Performance by Strategy} oraz tabela tagów}

% --- 5. Kalendarz ---
\subsection*{Kalendarz tradingowy}
Widok kalendarza łączy w~jednej siatce miesięcznej dwa źródła danych: własne transakcje użytkownika oraz wydarzenia makroekonomiczne pobierane co~6~godzin z~ForexFactory (klasa \texttt{MacroEventService} z~mechanizmem cache i~fallbackiem). Każda komórka dnia może wyświetlać: numer dnia, kropki impactu (czerwone, pomarańczowe, żółte, odpowiadające High/Medium/Low), miniaturową informację o~liczbie transakcji w~danym dniu wraz z~sumą P\&L (kolorowane) oraz cienki kolorowy pasek w~stopce komórki sygnalizujący wynik dnia (zielony~=~zysk, czerwony~=~strata, szary~=~break-even). Nad~kalendarzem znajduje się pasek filtrów krajów (USD, EUR, GBP, ...) pozwalający ograniczyć wyświetlane eventy do~rynków interesujących tradera. Klik w~komórkę dnia otwiera modal z~pełną listą transakcji i~eventów makro, z~dodatkowymi metrykami dnia (W/L, P\&L, ilość zamkniętych pozycji).
\screenshotPlaceholder{calendar-month}{Kalendarz tradingowy, widok miesięczny z~transakcjami i~eventami makro}
\screenshotPlaceholderTall{calendar-day}{Kalendarz, modal detali dnia z~listą transakcji i~wydarzeń makroekonomicznych}

% --- 6. Playbook ---
\subsection*{Moduł Playbook, biblioteka strategii}
Playbook jest osobistym katalogiem strategii tradingowych użytkownika. Lista strategii prezentowana jest jako siatka kart z~obrazkiem podglądu, tytułem, krótkim opisem, tagami (kolorowane chipy) i~wycinkiem treści. Klik w~kartę otwiera widok \emph{read-only} z~wyrenderowanym Markdownem i~zaznaczoną checklistą; klik w~przycisk \emph{Edit} przełącza do~edytora dwupanelowego: po~lewej pole tekstowe Markdown (czcionka monospace), po~prawej tab \emph{Preview} z~live-podglądem wyrenderowanym przez \texttt{react-markdown} i~\texttt{remark-gfm}. Edytor pozwala na~upload obrazka ilustrującego setup (drag~\&~drop), zarządzanie checklistą pre-trade (lista pozycji, dodawanie po~Enter, usuwanie krzyżykiem) oraz tagi rozdzielone przecinkami. Każda strategia może być wybrana w~modalu transakcji, co~powoduje automatyczne wczytanie jej checklisty do~formularza tradingowego.
\screenshotPlaceholder{playbook-list}{Playbook, lista strategii w~układzie kart z~tagami i~obrazkami podglądu}
\screenshotPlaceholderTall{playbook-view}{Playbook, widok read-only strategii z~wyrenderowanym Markdownem}

% --- 7. Long Term ---
\subsection*{Moduł Long~Term, planowanie strategiczne}
Long~Term to~moduł wykraczający poza horyzont dnia, który organizuje myślenie inwestora w~cztery kategorie: \textbf{THEME} (motywy rynkowe wymagające obserwacji przez tygodnie/miesiące), \textbf{WATCHLIST} (instrumenty oczekujące na~trigger lub potwierdzenie), \textbf{GOAL} (cele kwartalne procesowe i~portfelowe) oraz \textbf{REVIEW} (okresowe przeglądy weryfikujące, co~jest jeszcze aktualne). Każdy wpis ma~swój status (\texttt{ACTIVE}, \texttt{ON\_WATCH}, \texttt{BUILDING}, \texttt{COMPLETED}, \texttt{INVALIDATED}) prezentowany jako kolorowy chip oraz osobne sekcje formularza: \emph{Title}, \emph{Asset}, \emph{Horizon} (1-3M, 3-6M, ...), \emph{Thesis} (teza), \emph{Trigger~Plan} (warunek wejścia), \emph{Invalidation} (warunek unieważnienia tezy) i~\emph{Notes}. U~góry strony cztery karty statystyczne agregują liczbę aktywnych wpisów per kategoria, a~pasek filtrów pozwala ograniczyć widok do~jednej kategorii lub do~wyników wyszukiwania pełnotekstowego we~wszystkich polach.
\screenshotPlaceholder{long-term-list}{Long Term, lista wpisów pogrupowanych w~kategoriach THEME, WATCHLIST, GOAL, REVIEW}

% --- 8. Chart ---
\subsection*{Moduł Chart, integracja z~TradingView}
Strona Chart jest pełnoekranowym widokiem zawierającym osadzony widget TradingView (skrypt \texttt{tv.js}). Daje on~tradera dostęp do~pełnej analizy technicznej bezpośrednio w~aplikacji: świece, wskaźniki techniczne (RSI, MACD, średnie kroczące, Bollinger Bands), narzędzia rysowania (linie trendu, Fibonacci, fale Elliotta), wybór timeframe'u (1m -- 1M), zmianę symbolu (akcje, forex, crypto, indeksy) oraz wbudowany kalendarz wydarzeń. Widget jest skonfigurowany w~motywie ciemnym pasującym do~reszty aplikacji, z~domyślną parą EUR/USD na~OANDA. Inicjalizacja oraz cleanup widgetu odbywa się w~hooku \texttt{useEffect} (zapobiega wyciekom DOM przy odmontowaniu komponentu).
\screenshotPlaceholder{chart}{Moduł Chart, widget TradingView z~analizą techniczną pary EUR/USD}

% --- 9. Admin ---
\subsection*{Panel administracyjny}
Panel administracyjny dostępny jest wyłącznie dla~użytkowników z~rolą \texttt{ADMIN} (warunek wyświetlenia linku w~sidebarze --- na~podstawie pola \texttt{role} z~payloadu JWT --- oraz warunek autoryzacyjny po~stronie backendu w~\texttt{SecurityConfig}). Panel składa się z~trzech zakładek: \textbf{Users} (tabela wszystkich kont z~kolumnami: username, email, role, liczba transakcji, liczba wpisów dziennika, data ostatniej transakcji, status active/disabled oraz akcje: \emph{Disable/Enable}, \emph{Reset~PW} i~\emph{Delete}), \textbf{Statistics} (siatka kart z~metrykami systemowymi: \emph{Total~Users}, \emph{Active~Users~30d}, \emph{Total~Trades}, \emph{Journal~Entries}, \emph{Avg~Trades/User}) oraz \textbf{Config} (edytor par klucz-wartość konfiguracji systemowej, np. limity wielkości pliku, flagi feature toggle). Reset hasła otwiera osobny modal z~polem nowego hasła (wymagane min. 8~znaków) i~przyciskiem \emph{Show/Hide}. Każda destrukcyjna akcja (Delete, Reset~PW) wymaga potwierdzenia w~globalnym \texttt{ConfirmDialog}.
\screenshotPlaceholder{admin-users}{Panel administracyjny, zakładka Users z~listą kont i~akcjami}

% ============================================================
\section{Wnioski z~ćwiczenia}
% ============================================================

Wszystkie zakładane funkcjonalności zostały zrealizowane w~stopniu zadowalającym. Aplikacja umożliwia pełne zarządzanie transakcjami z~uwzględnieniem parametrów psychologicznych, dokumentowanie strategii w~module Playbook, prowadzenie dziennika psychologicznego, analizę wyników za pomocą zaawansowanych metryk i~wizualizacji oraz planowanie długoterminowe.

Szczególnie dobrze udało się zrealizować:
\begin{itemize}
  \item \textbf{Architekturę warstwową backendu}, podział na kontrolery, serwisy, repozytoria i~mappery zapewnił czytelność kodu i~łatwość rozbudowy. System 13 migracji Flyway zagwarantował powtarzalność schematów bazy danych.
  \item \textbf{System bezpieczeństwa}, autentyfikacja JWT, hashowanie BCrypt, ochrona przed enumeracją emaili, weryfikacja przynależności zasobów i~konfiguracja CORS tworzą spójną warstwę ochronną.
  \item \textbf{Optymalizację frontendu}, lazy loading stron, memoizacja obliczeń (\texttt{useMemo}), mechanizm auto-retry i~globalne konteksty (Toast, Confirm) zapewniają płynne działanie interfejsu.
  \item \textbf{Konteneryzację}, Docker Compose umożliwia uruchomienie całego systemu jednym poleceniem, co znacznie upraszcza zarówno lokalne testowanie, jak i~wdrożenie produkcyjne.
\end{itemize}

Elementy, które można by usprawnić w~przyszłości:
\begin{itemize}
  \item Dodanie testów jednostkowych i~integracyjnych (JUnit, React Testing Library) w~celu zwiększenia niezawodności.
  \item Wdrożenie WebSocket do aktualizacji cen rynkowych w~czasie rzeczywistym zamiast odpytywania REST.
  \item Rozbudowa modułu Analytics o~analizę korelacji między nastrojem z~dziennika a~wynikami transakcji.
  \item Dodanie eksportu danych do CSV/PDF na potrzeby raportowania.
\end{itemize}

Projekt pozwolił na praktyczne zastosowanie wzorców architektonicznych (MVC, Builder, DTO, Repository), poznanie mechanizmów bezpieczeństwa webowego oraz doświadczenie pełnego cyklu tworzenia aplikacji, od projektowania bazy danych, przez implementację logiki biznesowej i~interfejsu, po konteneryzację i~wdrożenie na platformę chmurową.

\end{document}
